\section*{Introduction}
\addcontentsline{toc}{section}{Introduction}

\textbf{The problem.} Fix a single polynomial equation $P(a, u) = 0$ with integer coefficients, of total degree at most two in the parameters $a ∈ ℤ^{k}$ and the unknowns $u ∈ ℤ^{r}$ together, and ask which relations on ℤᵏ it can define by existential projection: $a ∈ A$ if and only if some integer vector u satisfies $P(a, u) = 0$. The class of such relations, over all r and all P of total degree at most two, is the single-quadratic existential fragment SE₂(k). Nothing restricts the number of unknowns, so the fragment is not small: it contains every non-semilinear set that a single quadratic can carve out, and its fibres already include Pell-type infinite orbits. The question studied here is how much this one degree bound constrains the fragment, and where the fragment sits between the classes it is squeezed between.

\textbf{The boundary it lies on.} Two classical results frame the question. Solvability of a single quadratic equation in integers is decidable (Grunewald–Segal \cite{GS04}), so a single quadratic atom cannot encode an undecidable relation. On the other side, every recursively enumerable relation is Diophantine (Davis–Putnam–Robinson–Matiyasevich \cite{DPR61,Mat70}), and Skolem's degree reduction \cite{Sko38} turns any polynomial equation into a system of quadratics, or into a single quartic. So the passage from \emph{one} quadratic equation to a primitive-positive system of quadratics — equivalently, to a single equation of degree four — crosses from a decidable class to the whole of the recursively enumerable relations. The fragment SE₂ sits on that boundary, and this paper is about what happens there: although the fragment already supports non-semilinear behaviour and Pell-type infinite orbits, its finite fibres and total functions satisfy uniform polynomial restrictions, and total degree and formula shape jointly control fibre structure, functional growth, closure under composition, and computability.

\textbf{Main result: the functions one quadratic equation defines.} A total function $f : ℤ^{n} → ℤ^{m}$ is SE₂-definable when its graph is. Say that $f$ has a \emph{common slope} if there are $M ≥ 1$ and an integer $m × n$ matrix $B$ with $f(a + Mt) = f(a) + Bt$ for all $a, t ∈ ℤ^{n}$: one rational slope $B/M$ on every residue class modulo $M$. Assume $\mathrm{TCER}$, an eventual representation property of positive definite ternary lattice cosets that follows from published work of Kane and Kim (\cref{lemma:10.8}). Then a function with two or more outputs is SE₂-definable exactly when it has a common slope (\cref{proposition:10.9}). A function $f : ℤ^{n} → ℤ$ is SE₂-definable exactly when it has a common slope, or agrees on $ℤ^{n}$ with a polynomial with rational coefficients of total degree at most two, or there are integer vectors $p, q$ and integers $c, c'$, with $p - q$ having even coordinates, $p$ having an odd coordinate and $c - c'$ odd, such that $2f(a) = p · a + c$ when $p · a + c$ is even and $2f(a) = q · a + c'$ otherwise (\cref{proposition:10.14}). The three families overlap, and none of them can be dropped. The converse constructions use no analytic input, and neither does the case without unknowns, where exactly the last two families occur (\cref{proposition:10.16}). The classification is proved in \cref{sec:10}. The sharp bound $n ≤ m$ for SE₂-definable injections (\cref{theorem:4.1}) keeps its own proof, which uses no analytic input.

\textbf{Finite fibres.} The first group of results says that a single quadratic cannot spread a finite fibre far. For every P of total degree at most two there is one constant C, depending on P alone, such that whenever the fibre over a parameter point a is finite, every element of it has sup-norm at most C·N², where N bounds $‖a‖_{∞}$ (\cref{theorem:2.1}), and the fibre has at most (2CN² + 1)ᵐ elements (\cref{theorem:2.2}) — the same constant serving both bounds. The uniformity in a is the point: the bound is not a search bound for the equation in front of us but a single constant valid at every parameter point, which is what the rest of the paper rests on.

\textbf{Definable functions.} A total function f : $ℤ^{n} → ℤ^{m}$ is SE₂-definable when its graph is. The finite-fibre bound forces such functions to grow at most quadratically (\cref{theorem:3.2}); they are computable (\cref{theorem:3.3}); and they are not closed under composition, since $x ↦ x^{2}$ is definable while $x ↦ x^{4}$ is not (\cref{theorem:3.4}). An SE₂-definable injection $ℤ^{n} → ℤ^{m}$ must satisfy $n ≤ m$ (\cref{theorem:4.1}), and conversely one exists whenever $n ≤ m$ (\cref{theorem:4.3}), so no single quadratic defines the graph of an injection $ℤ^{n} → ℤ$ for $n ≥ 2$ and the one-sorted coding arity is exactly one; and an integer polynomial has SE₂-definable graph exactly when its total degree is at most two (\cref{theorem:5.1}). The growth bound and the non-closure follow from the same constant C; that constant alone gives only $n ≤ 2m$ for injections (\cref{remark:4.5}), and the sharp bound $n ≤ m$ comes from the structural alternative of \cref{theorem:4.2} — a box count on the linear-growth side and a collision $f(ρ + Mv) = f(ρ − Mv)$ on the coset side.

\textbf{Normal forms.} Beyond growth, the equation fixes the shape of a definable function. \cref{theorem:4.2} gives every SE₂-definable total function linear growth or a weak vector coset normal form. The full form --- one modulus and, on every residue class, a vector quadratic polynomial --- holds unconditionally except when the quadratic part in the unknowns is semidefinite of polar rank at least three, and everywhere under one analytic input, $\mathrm{TCER}$, the eventual representation of large integers by an individual positive definite ternary lattice coset (\cref{proposition:10.6,proposition:10.7}). $\mathrm{TCER}$ is a consequence of published work on theta series of lattice cosets (\cref{lemma:10.8}) and enters the formal development as an explicit hypothesis. With at least two outputs the branch proofs give the common slope of the main result. Consequently $a ↦ (a^{2}, 0)$ is not SE₂-definable under $\mathrm{TCER}$ although $a ↦ a^{2}$ is (\cref{proposition:10.10}), while a full form alone does not yield a common slope (\cref{proposition:10.11}). Appending a zero output keeps a function SE₂-definable exactly when it has a common slope (\cref{proposition:10.15}).

\textbf{Infinite fibres: finite or structured.} The second group describes the infinite fibres positively. \cref{theorem:8.5} asserts, unconditionally and with one constant C for all parameter points, that every scalar fibre of every equation of total degree at most two is either bounded by C(1 + $‖a‖_{∞}$)² or carries a structured certificate: arithmetic data of one of three kinds — a polynomial progression, a Pell-unit orbit, or a moving-modulus step orbit — that generates an infinite subset of the fibre and is represented in the formalization by a typed realization whose fields witness an infinite arithmetic output subset. The certificate pins an infinite output subset, not the whole fibre, and the three kinds are a sufficient language rather than a classification. The engine behind both groups of results is an arithmetic projection dichotomy for homogeneous quadratic forms (\cref{theorem:9.2,theorem:9.3}): for a fixed form, the set of natural numbers c admitting an integer witness of an affine perturbation is either infinite or bounded by a constant depending on the form alone.

\textbf{The computability phase diagram.} Combining the structural results with the two classical inputs places the fragment exactly. Under code-uniform decidability of quadratic solvability and DPRM over ℕ,

\[
SE_{2}(k) ⊊ \mathrm{REC}(k) ⊊ \mathrm{RE}(k) = PP_{2}(k) = SE_{4}(k) \quad\text{for every } k ≥ 1,
\]

and, for total functions, $\mathrm{Fun}(SE_{2}) ⊊ \mathrm{TotComp} = \mathrm{Fun}(PP_{2}) = \mathrm{Fun}(SE_{4})$ at every positive input and output arity (\cref{theorem:6.8,theorem:7.1}). At arity zero all five classes coincide, which is why every strictness statement carries an arity hypothesis. Two of the arrows are proved without either input: $\mathrm{REC} ⊊ \mathrm{RE}$ (\cref{theorem:6.1}) and $\mathrm{Fun}(SE_{2}) ⊊ \mathrm{TotComp}$ (\cref{theorem:3.5}); the inclusion $SE_{2} ⊆ \mathrm{REC}$ uses the quadratic-solvability principle, and the equalities with PP₂ and SE₄ use DPRM. The strictness $SE_{2} ⊊ \mathrm{REC}$ at arity one is the delicate arrow: it needs the uniform, code-level form of decidability, and is obtained by an effective diagonalization over a coded family (\cref{theorem:6.5}, \cref{corollary:6.6}).

\textbf{What is new and what is used.} We separate the contributions by kind, so that a reader can see at once what this paper claims and what it takes from the literature.

\begin{itemize}
\item \emph{Core mathematical novelty candidates.} The uniform finite-fibre height and cardinality bounds with a single constant (\cref{theorem:2.1,theorem:2.2}); quadratic growth of SE₂-definable functions (\cref{theorem:3.2}) and their non-closure under composition as a structural consequence (\cref{theorem:3.4}); the sharp injective dimension theorem $n ≤ m$ with its converse (\cref{theorem:4.1,theorem:4.3}) and the structural alternative behind it (\cref{theorem:4.2}); the graph classification $\mathrm{graph}(p) ∈ SE_{2} ⟺ \mathrm{deg}\ p ≤ 2$ (\cref{theorem:5.1}); the unconditional finite-or-structured theorem (\cref{theorem:8.5}); and the arithmetic projection dichotomy in bounded and data-carrying form for every number of unknowns (\cref{theorem:9.2,theorem:9.3}). We call them candidates because our literature search (\cref{sec:11}) found no statement containing them, not because novelty has been externally adjudicated.
\item \emph{Definable functions.} Under $\mathrm{TCER}$, the exact characterization of the SE₂-definable total functions: a common slope for two or more outputs (\cref{proposition:10.9}), three families for one output (\cref{proposition:10.14}), with converse constructions that need no analytic input; the case without unknowns, unconditionally (\cref{proposition:10.16}); the padding criterion (\cref{proposition:10.15}) and presentations with unique unknowns (\cref{proposition:10.17}); and the reduction of the full vector coset normal form to the nonzero semidefinite case (\cref{proposition:10.6}), within which polar rank at most two is settled without further input, with the full form under $\mathrm{TCER}$ (\cref{proposition:10.7}). The components are known or elementary: the one-hot quadratic penalty behind the converse, the rigidity of square-valued quadratic forms (\cref{lemma:10.4}) and the two-coset lemma (\cref{lemma:10.13}); what is added is the characterization of the functions, compared with Presburger-definable functions in \cref{sec:11}. $\mathrm{TCER}$ is taken from published work (\cref{lemma:10.8}).
\item \emph{Standard logical consequences and new synthesis.} The computability of SE₂-definable total functions from a total single-valued c.e. graph (\cref{theorem:3.3}) and the exactness of $\mathrm{Fun}(SE_{2}) ⊊ \mathrm{TotComp}$ by arity (\cref{theorem:3.5}); the effective diagonalization over a coded family and its arity-one consequence (\cref{theorem:6.5,corollary:6.6}), a standard argument made effective here by the code bridge; Skolem's degree reduction $PP_{2} = SE_{4} = \mathrm{Dioph}$ (\cref{theorem:6.2}), the transport of DPRM from ℕ to ℤ (\cref{theorem:6.7}), the code bridge from our finite quadratic codes to the coefficient-table input of \cite{GS04} (\cref{theorem:6.4}), and the exact relation and function phase diagrams (\cref{theorem:6.8,theorem:7.1}). These are known arrows or standard arguments assembled into one exact picture; the arrows are not claimed as new, the exactness of the picture is.
\item \emph{Classical inputs.} Decidability of integer quadratic solvability in its code-uniform form (\cref{definition:1.5}, after \cite{GS04}) and DPRM over ℕ (\cref{definition:1.6}), the theorem of \cite{DPR61,Mat70}. Both are published results; they enter as explicit named hypotheses of the theorems that use them, and neither is an axiom of the formalization.
\item \emph{Formalization contributions.} The typed certificate language with per-kind soundness (\cref{definition:8.2,theorem:8.3}), the quadratic-code and DPRM transport layers, and the Lean development itself: Sections 2–10 and \cref{sec:C} are Lean 4 declarations checked against mathlib, \cref{sec:10} with $\mathrm{TCER}$ as an explicit hypothesis and apart from \cref{lemma:10.8}; every statement in the statement-fidelity registry is frozen and diffed against the development, every root's axiom footprint is within the standard three, and the artifact ships the checks that recompute these facts (\cref{sec:12}).
\end{itemize}

\textbf{Proof strategy.} Everything on the finite side flows from one decomposition and one dichotomy. A polynomial of total degree at most two in (a, u) is H(c, u) + Σ ℓᵢ(a)·xᵢ + q(a), with H the homogeneous quadratic component in the output-and-witness variables and only the affine part depending on the parameter (\cref{sec:2}). The arithmetic projection dichotomy (\cref{sec:9}, proved in full in \cref{sec:F}) says that for a fixed H the set of c admitting a witness for any affine perturbation is either infinite or bounded by $K·((1 + Σ|ℓ_{i}|)^{2} + |q|)$ — a constant $K$ depending on $H$ alone, times a quadratic height in the affine data $(ℓ, q)$ (\cref{definition:9.1}); the constant $K$ is not effective, and \cref{sec:12} and the trust document say so. Applied to the fibre equation this gives \cref{theorem:2.1} at once, and the growth and non-closure consequences for functions follow by reading the graph of f as a fibre. The sharp injective bound $n ≤ m$ of \cref{theorem:4.1} does not follow from the fibre bound alone: it needs the structural alternative of \cref{theorem:4.2}, a box count on the linear-growth side and a collision $f(ρ + Mv) = f(ρ − Mv)$ on the coset side. The infinite side is the same dichotomy in data-carrying form: where the projection is infinite, the proof produces the certificate rather than merely asserting infinitude, and the three kinds capture three recurring arithmetic mechanisms emitted across the proof branches (\cref{remark:8.8}). The computability side is diagonalization over an explicitly coded family, made effective by the code bridge.

\textbf{Organization.} \cref{sec:1} fixes the fragments SE₂, REC, RE, PP₂, SE₄ and the two external inputs. Sections 2–5 prove the finite-fibre bounds and their consequences for functions, injections and polynomial graphs. Sections 6 and 7 assemble the relation and function phase diagrams. \cref{sec:8} introduces structured certificates and proves the finite-or-structured theorem; \cref{sec:9} states the arithmetic projection dichotomy that drives it. \cref{sec:10} proves the normal forms and the classification of definable functions. \cref{sec:11} places the results against the literature, and \cref{sec:12} describes the formalization and what it does and does not certify. \cref{sec:A} collects open problems; Appendices B–D collect capacity refinements, the multi-parameter universal-pair question that motivated the fragment, and technical lemmas; \cref{sec:E} is the statement map from paper labels to Lean names; \cref{sec:F} is the complete proof of the projection dichotomy.

